    \documentclass[12pt,a4paper]{article}
    \usepackage[a4paper,top=2.5cm,bottom=2.5cm,left=2.5cm,right=2.5cm]{geometry}
    \usepackage[T5]{fontenc}
    \usepackage[utf8]{inputenc}
    \usepackage[english,vietnamese]{babel}
    \usepackage{amsmath,amssymb}
    \usepackage{graphicx}
    \usepackage{pgfplots}
    \pgfplotsset{compat=1.18}
    \usepackage{booktabs}
    \usepackage{array}
    \usepackage{float}
    \usepackage{listings}
    \usepackage{hyperref}
\usetikzlibrary{positioning,arrows.meta,shapes.geometric,fit,babel}
\tikzset{
    box/.style={draw, rounded corners=2pt, align=center, font=\small, inner sep=4pt, minimum height=8mm},
    decision/.style={draw, diamond, aspect=2.2, align=center, font=\small, inner sep=2pt},
    terminal/.style={draw, rounded corners=6pt, align=center, font=\small\bfseries, inner sep=4pt, minimum height=7mm},
    arr/.style={-{Stealth[length=2mm]}, thick},
    lbl/.style={font=\footnotesize, inner sep=3pt},
}

    \hypersetup{
        colorlinks=true,
        linkcolor=blue,
        filecolor=magenta,
        urlcolor=cyan,
        hypertexnames=false,
    }

    \setlength{\emergencystretch}{2em}

    \title{Báo cáo Đồ án 2: Griductive Solver}

    \begin{document}

    \begin{titlepage}

    \newcommand{\HRule}{\rule{\linewidth}{0.5mm}}

    \centering
    \vspace*{0.5cm}

    \textsc{\LARGE Trường Đại học Khoa học Tự nhiên}\\[0.35cm]
    \textsc{\Large Ngành Trí tuệ Nhân tạo}\\[0.35cm]
    \textsc{\large Môn học: Cơ sở Trí tuệ Nhân tạo}\\[1.5cm]

    \includegraphics[width=0.30\textwidth]{image/logo-khtn.png}\\[1.5cm]

    \HRule \\[0.4cm]
    {\huge \bfseries Báo cáo Đồ án 2}\\[0.3cm]
    {\Large \bfseries Griductive Solver}\\[0.4cm]
    \HRule \\[2.0cm]

    {\renewcommand{\arraystretch}{1.25}%
    \noindent
    \begin{minipage}[t]{0.56\textwidth}
    \begin{flushleft}
    {\large \emph{Sinh viên thực hiện:}}\\[0.25cm]
    \normalsize
    \begin{tabular}[t]{@{}l@{\hspace{0.5em}--\hspace{0.5em}}l@{}}
    Nguyễn Huỳnh Hải Long & 25122002 \\
    Ngô Quang Minh        & 25122011 \\
    Võ Tấn Lợi            & 25122028 \\
    Hồ Trần Quốc Thắng    & 25122041 \\
    Nguyễn Lê Khánh Toàn  & 25122045
    \end{tabular}
    \end{flushleft}
    \end{minipage}\hfill
    \begin{minipage}[t]{0.42\textwidth}
    \begin{flushright}
    {\large \emph{Giáo viên hướng dẫn:}}\\[0.25cm]
    \normalsize
    \begin{tabular}[t]{@{}r@{}}
    NGND, GS. TS. Lê Hoài Bắc \\
    ThS. Nguyễn Ngọc Đức
    \end{tabular}
    \end{flushright}
    \end{minipage}}

    \vfill

    \end{titlepage}

    \tableofcontents
    \newpage

    % ==================== PHẦN 1 ====================
    \section{Giới thiệu và Tổng quan bài toán}

    \subsection{Giới thiệu trò chơi Griductive}
    Griductive \cite{griductive_game} là trò chơi giải đố suy luận trên lưới ô vuông kích thước $N \times N$ ($N \in \{3, 4, 5\}$), trong đó mỗi ô chứa một nhân vật mang trạng thái ẩn là \emph{Criminal} (tội phạm) hoặc \emph{Innocent} (người vô tội). Mục tiêu của người chơi là dựa vào các manh mối trên những thẻ đã lật ngửa để lần lượt xác định danh tính của toàn bộ nhân vật còn lại. Điểm khó của trò chơi là người chơi \emph{không được đoán}: mỗi phán quyết chỉ được chấp nhận khi trạng thái đó là hệ quả logic tất yếu từ các dữ kiện hiện có; nếu phán quyết mang tính phỏng đoán hoặc mâu thuẫn, thẻ sẽ không được lật.

    Bài toán phù hợp để áp dụng các phương pháp suy diễn tự động dựa trên logic mệnh đề, mã hóa dạng chuẩn hội (CNF) và bộ giải thỏa mãn Boolean (DPLL SAT solver). Đặc biệt, vì mọi manh mối đã lật đều luôn đúng và bộ giải chỉ được tiếp cận trạng thái công khai theo từng bước, đây là môi trường tiêu biểu để thiết kế và đánh giá tác tử suy luận logic chặt chẽ \cite{russell2020aima}.

    \subsection{Luật chơi và Các thành phần trong màn chơi}
    \subsubsection{Cấu trúc bàn chơi và Manh mối}
    \begin{itemize}
        \item \textbf{Cấu trúc lưới và Nhân vật:} Bàn chơi gồm $N$ hàng đánh số từ trên xuống ($1 \dots N$) và $N$ cột đánh chữ cái từ trái sang phải ($A, B, \dots$; ví dụ \texttt{C2}). Mỗi ô chứa một nhân vật với tên (xếp theo bảng chữ cái) và nghề nghiệp; các thông tin này không mang ý nghĩa logic trừ khi manh mối nhắc trực tiếp đến.

        \item \textbf{Trạng thái thẻ bài và Manh mối:} Ban đầu một số thẻ được lật ngửa để cung cấp dữ kiện khởi đầu, các thẻ còn lại úp mặt. Mọi manh mối đã lật đều đúng bất kể chủ nhân thẻ là \emph{Criminal} hay \emph{Innocent}, và được biểu diễn theo các mẫu có cấu trúc chuẩn.

        \item \textbf{Cơ chế suy diễn và Điều kiện thắng:} Ở mỗi lượt, người chơi chọn một nhân vật chưa xác định và gửi phán quyết \texttt{CRIMINAL} hoặc \texttt{INNOCENT}. Nếu phán quyết đúng và chứng minh được, thẻ sẽ lật ngửa để tiết lộ manh mối mới; người chơi thắng khi mở chính xác toàn bộ $N \times N$ ô trên lưới.
    \end{itemize}

    \subsubsection{Các vùng manh mối và Cơ chế phản hồi}
    \begin{itemize}
        \item \textbf{Chỉ kết luận khi chứng minh được (No-guess rule):} Mỗi màn chơi hợp lệ có duy nhất một phép gán hoàn chỉnh thỏa mãn toàn bộ manh mối, và người chơi chỉ được phép mở thẻ khi trạng thái nhân vật bị ép buộc logic bởi tập manh mối hiện tại.

        \item \textbf{Các vùng manh mối (Clue regions):} Một manh mối có thể tham chiếu đến một trong các vùng trên lưới ô vuông:
        \begin{itemize}
            \item \emph{Hàng và Cột (Row \& Column):} Tập hợp tất cả các ô có cùng chỉ số hàng hoặc cùng chữ cái cột.
            \item \emph{Lân cận 8 hướng (Neighbors):} Các ô tiếp xúc ngang, dọc hoặc chéo với một ô cho trước (tối đa 8 ô, không tính chính ô đó).
            \item \emph{Danh sách ô chỉ định (Explicit list):} Tập hợp tường minh các tọa độ ô phân biệt được liệt kê trong manh mối.
        \end{itemize}

        \item \textbf{Cơ chế phản hồi phán quyết:} Khi người chơi đưa ra phán quyết chưa hợp lệ, hệ thống giữ nguyên bàn chơi mà không lật thẻ ẩn, đồng thời phân biệt rõ giữa \texttt{NOT\_PROVABLE} (chưa đủ dữ kiện để chứng minh cả hai chiều) và \texttt{CONTRADICTED} (mâu thuẫn với các manh mối đã biết).
    \end{itemize}
    % ==================== PHẦN 2 ====================
    \section{Phân công công việc}

\begin{table}[H]
\centering
\renewcommand{\arraystretch}{1.3}
\begin{tabular*}{\textwidth}{@{\extracolsep{\fill}}clp{0.47\textwidth}c@{}}
\toprule
\textbf{STT} & \textbf{Họ và tên} & \textbf{Công việc phụ trách} & \textbf{Hoàn thành} \\
\midrule
1 & Nguyễn Huỳnh Hải Long & Xây dựng bộ mã hóa CNF (\texttt{encoder.py}): ánh xạ biến tất định, mã hóa sáu mẫu manh mối cốt lõi và hai manh mối mở rộng \texttt{PARITY}, \texttt{INTERSECTION}; phối hợp chốt các kiểu dữ liệu chung; viết kiểm thử cho bộ mã hóa. & 100\% \\
2 & Ngô Quang Minh        & Xây dựng Game Engine và bộ đánh giá ngữ nghĩa trực tiếp, viết kiểm thử đối chiếu với bộ mã hóa CNF; xây dựng bộ chạy thực nghiệm và xuất CSV; viết phần Thực nghiệm; quay video demo. & 100\% \\
3 & Võ Tấn Lợi            & Xây dựng vòng lặp suy diễn (\texttt{agent/runner.py}) và vết suy diễn (\texttt{agent/trace.py}); cùng hoàn thiện lớp tác tử; tạo bộ dữ liệu màn chơi ban đầu. & 100\% \\
4 & Hồ Trần Quốc Thắng    & Thiết lập CI; thiết kế kiểu dữ liệu chung, hệ tọa độ và vùng manh mối; cài đặt bộ giải DPLL, bộ nạp JSON và bộ kiểm định tính duy nhất; thiết kế lại bộ màn chơi \texttt{puzzles/}; xây dựng giao diện Pygame và \texttt{main.py}. & 100\% \\
5 & Nguyễn Lê Khánh Toàn  & Xây dựng tác tử suy luận (\texttt{agent/inference.py}): dựng $KB_t$ từ trạng thái công khai và phân loại hệ quả logic; dựng khung báo cáo. & 100\% \\
\bottomrule
\end{tabular*}
\caption{Bảng phân công công việc và mức độ hoàn thành của các thành viên.}
\label{tab:task_distribution}
\end{table}

    % ==================== PHẦN 3 ====================
    \section{Thiết kế hệ thống}

    \subsection{Kiến trúc mã nguồn}
Mã nguồn nằm trong gói \texttt{griductive/} và tách hai vai trò theo yêu cầu của đề (Hình~\ref{fig:architecture}). \texttt{GridGameEngine} giữ toàn bộ màn chơi, gồm đáp án và manh mối của các thẻ úp, và chỉ nhận màn đã qua bộ kiểm định. Tác tử chỉ nhận \texttt{PublicState}, gồm bàn chơi, các phán quyết đã xác nhận và manh mối của thẻ đang mở, rồi dựng $KB_t$ từ đó.

Khi nhận một phán quyết, engine phân loại ô được chọn trên chính $KB_t$ công khai. Engine trả về \texttt{NOT\_PROVABLE} nếu cả hai trạng thái còn khả dĩ, và \texttt{CONTRADICTED} nếu trạng thái ngược lại là hệ quả logic. Nếu phán quyết là hệ quả logic, engine đối chiếu với đáp án rồi công bố trạng thái và manh mối của ô trong cùng một lần cập nhật. Phán quyết bị từ chối không làm thay đổi bàn chơi; \texttt{Restart} phục hồi các ô mở ban đầu.

\begin{figure}[H]
\centering
\begin{tikzpicture}
\node[box] (ui) at (0,0) {GUI / Runner};
\node[box] (eng) at (8.6,0) {\textbf{Game Engine}\\ \footnotesize giữ đáp án và manh mối thẻ úp};
\node[box] (agent) at (0,-4.4) {\textbf{Logic Agent}};
\node[box] (enc) at (4.3,-4.4) {CNF Encoder};
\node[box] (sat) at (8.6,-4.4) {DPLL Solver};
\coordinate (lowest) at (0,-6.1);
\draw[arr] ([yshift=2.5mm]ui.east) -- node[lbl, above] {phán quyết (ô, trạng thái)} ([yshift=2.5mm]eng.west);
\draw[arr] ([yshift=-2.5mm]eng.west) -- node[lbl, below] {kết quả + manh mối mới} ([yshift=-2.5mm]ui.east);
\draw[arr] (ui) -- node[lbl, left, align=right] {\texttt{PublicState}\\(Hint, Auto solve)} (agent);
\draw[arr] (eng.south) -- ++(0,-1.4) -| node[lbl, above, pos=0.25] {\texttt{PublicState} (chấm phán quyết)} ([xshift=7mm]agent.north);
\draw[arr] (agent) -- node[lbl, above, align=center] {manh mối,\\phán quyết} (enc);
\draw[arr] (enc) -- node[lbl, above] {$KB_t$} (sat);
\draw[arr] (sat.south) -- ++(0,-0.8) -| node[lbl, below, pos=0.25] {SAT / UNSAT cho từng giả thiết} (agent.south);
\node[draw, dashed, rounded corners, inner sep=5mm, fit=(agent)(enc)(sat)(lowest)] (pub) {};
\node[lbl, anchor=south east] at (pub.north east) {chỉ dùng thông tin công khai};
\end{tikzpicture}
\caption{Tách vai trò: chỉ Game Engine giữ dữ liệu bí mật; tác tử, bộ mã hóa và bộ giải SAT chỉ làm việc trên \texttt{PublicState}.}
\label{fig:architecture}
\end{figure}

    \subsection{Định dạng và bộ màn chơi}
Mỗi màn là một file JSON gồm sáu khóa: \texttt{id}, \texttt{size}, \texttt{characters}, \texttt{clues}, \texttt{solution} và \texttt{initially\_revealed}. Tọa độ được ghi dạng \texttt{A1}, \texttt{C2} và được đổi sang \texttt{Cell(row, col)} đánh số từ $0$, sắp theo thứ tự row-major. Bảng~\ref{tab:json_schema} liệt kê các trường của từng loại manh mối và vùng. \texttt{JsonPuzzleLoader} kiểm tra chặt cấu trúc và báo lỗi kèm vị trí (ví dụ \texttt{clues[2].k: must be in [0, 3], got 4}), không tự sửa dữ liệu. Bộ nạp từ chối các file có khóa thiếu hoặc lạ, tọa độ ngoài bàn, tên nhân vật không theo thứ tự chữ cái, \texttt{id} trùng, hoặc $k$ nằm ngoài $[0, |R|]$. Các kiểm tra ngữ nghĩa (đáp án thỏa mọi manh mối, nghiệm duy nhất) do bộ kiểm định ở Mục~\ref{sec:uniqueness} đảm nhận.

\begin{table}[H]
\centering\small
\renewcommand{\arraystretch}{1.2}
\begin{tabular}{@{}ll@{}}
\toprule
\textbf{Loại manh mối} & \textbf{Trường riêng (ngoài \texttt{id}, \texttt{owner}, \texttt{type})} \\
\midrule
\texttt{FACT} & \texttt{target}, \texttt{status} \\
\texttt{SAME}, \texttt{DIFFERENT} & \texttt{left}, \texttt{right} \\
\texttt{EXACTLY}, \texttt{AT\_LEAST}, \texttt{AT\_MOST} & \texttt{k}, \texttt{region} \\
\texttt{PARITY} & \texttt{parity} (\texttt{EVEN}/\texttt{ODD}), \texttt{region} \\
\texttt{INTERSECTION} & \texttt{k}, \texttt{regions} (mảng từ hai vùng trở lên) \\
\midrule
\multicolumn{2}{@{}l@{}}{Vùng: \texttt{\{"type":"ROW","row":2\}}, \texttt{\{"type":"COLUMN","column":"C"\}},} \\
\multicolumn{2}{@{}l@{}}{\phantom{Vùng: }\texttt{\{"type":"NEIGHBORS","cell":"B2"\}}, \texttt{\{"type":"EXPLICIT","cells":["A1","C2"]\}}} \\
\bottomrule
\end{tabular}
\caption{Lược đồ JSON của manh mối và vùng.}
\label{tab:json_schema}
\end{table}

Thư mục \texttt{puzzles/} gồm mười màn theo chủ đề Dragon Ball: ba màn cho mỗi kích thước $3\times3$, $4\times4$, $5\times5$, và màn \texttt{4x4-04} dành cho hai manh mối mở rộng. Mọi màn đều đọc được, có đáp án thỏa mọi manh mối, có nghiệm duy nhất và được tác tử giải trọn mà không phải đoán. Ở màn \texttt{4x4-04}, nếu loại riêng các manh mối \texttt{PARITY}, hoặc riêng các manh mối \texttt{INTERSECTION}, khỏi $KB_t$ thì tác tử đều bế tắc; như vậy cả hai loại manh mối mở rộng đều cần thiết để giải màn này.

\subsection{Bộ đánh giá ngữ nghĩa trực tiếp (Direct Semantic Evaluator)}
    \texttt{DirectClueEvaluator} nhận một manh mối, thông tin bàn chơi và một phép gán trạng thái đầy đủ. Hàm kiểm tra kích thước, tọa độ, kiểu trạng thái và tham số liên quan, sau đó đánh giá trực tiếp bằng phép so sánh hoặc đếm. Đầu vào không hợp lệ gây lỗi \texttt{ValueError}. Kết quả Boolean cho biết phép gán có thỏa mãn manh mối hay không.

    \begin{itemize}
        \item \texttt{FACT} so sánh trạng thái tại ô đích với trạng thái yêu cầu. \texttt{SAME} và \texttt{DIFFERENT} kiểm tra hai ô có cùng hay khác trạng thái.
        \item \texttt{EXACTLY}, \texttt{AT\_LEAST} và \texttt{AT\_MOST} đếm số ô \emph{Criminal} trong vùng rồi so sánh kết quả với ngưỡng $k$ theo thứ tự $=k$, $\ge k$ và $\le k$.
        \item Manh mối chẵn lẻ kiểm tra số ô \emph{Criminal} theo modulo $2$. Manh mối giao vùng lấy giao của các vùng trước khi kiểm tra số ô \emph{Criminal} có đúng bằng $k$ hay không.
    \end{itemize}

    Bộ đánh giá chạy độc lập với bộ mã hóa và bộ giải SAT. Kiểm thử đối chiếu duyệt cả 16 phép gán trên bàn $2 \times 2$ cho mỗi trường hợp manh mối được chọn. Với từng phép gán, kiểm thử cố định các biến chính trong CNF và kiểm tra có tồn tại giá trị cho các biến phụ để thỏa mãn CNF hay không. Kết quả này phải trùng với kết quả đánh giá trực tiếp. Các trường hợp được chọn bao gồm hai mẫu mở rộng và trường hợp giao vùng rỗng.

    \subsection{Giao diện đồ họa (GUI)}
Giao diện là ứng dụng desktop viết bằng Pygame~\cite{pygame}, khởi chạy bằng lệnh \texttt{python main.py}; ảnh nhân vật lấy từ Dragon Ball API~\cite{dragonball_api}. Lưới hiển thị nhãn hàng, cột, tên và nghề nghiệp của từng nhân vật; thẻ đã lật hiện thêm trạng thái và manh mối, còn thẻ úp không lộ thông tin nào. Người chơi chọn một thẻ rồi bấm \texttt{Innocent} hoặc \texttt{Criminal} (phím \texttt{I}/\texttt{C}). Nếu phán quyết được chấp nhận, thẻ lật lên và manh mối mới được thêm vào danh sách. Nếu bị từ chối, bàn chơi giữ nguyên và khung phản hồi phân biệt \texttt{NOT\_PROVABLE} với \texttt{CONTRADICTED} (Hình~\ref{fig:gui_reject}). Bấm vào một manh mối sẽ tô sáng các ô mà manh mối đó tham chiếu. Thanh điều khiển có bốn nút: \texttt{Load} chọn màn trong \texttt{puzzles/}, \texttt{Restart} đưa bàn về trạng thái ban đầu, \texttt{Hint} chỉ ra nhân vật chứng minh được tiếp theo cùng manh mối liên quan mà không đọc đáp án, và \texttt{Auto solve} cho tác tử giải từng bước. Khung \emph{Solver log} hiển thị số lượt gọi SAT, số lần chọn nhánh, lan truyền, quay lui, cùng số biến chính, biến phụ và số mệnh đề của $KB_t$. Giao diện chỉ đọc \texttt{PublicState} từ engine nên không thể làm lộ thẻ úp. Mọi lời gọi tới engine và tác tử đều chạy trên luồng nền để cửa sổ không bị treo.

\begin{figure}[H]
\centering
\includegraphics[width=0.9\linewidth]{image/gui-board.png}
\caption{Màn \texttt{4x4-01} sau khi nạp: hai thẻ mở ban đầu (\texttt{A1}, \texttt{D4}) và manh mối của Bardock đang được chọn, các ô lân cận của \texttt{A1} được viền xanh.}
\label{fig:gui_board}
\end{figure}

\begin{figure}[H]
\centering
\begin{minipage}[t]{0.49\linewidth}
\centering
\includegraphics[width=\linewidth]{image/gui-not-provable.png}
\end{minipage}\hfill
\begin{minipage}[t]{0.49\linewidth}
\centering
\includegraphics[width=\linewidth]{image/gui-contradicted.png}
\end{minipage}
\caption{Hai kiểu từ chối phán quyết. Trái: gọi Bulma (\texttt{C1}) là \emph{Criminal} khi chưa đủ dữ kiện, nhận \texttt{NOT\_PROVABLE}. Phải: gọi Broly (\texttt{B1}) là \emph{Innocent} trong khi các manh mối đã chứng minh Broly là \emph{Criminal}, nhận \texttt{CONTRADICTED}.}
\label{fig:gui_reject}
\end{figure}

% ==================== PHẦN 4 ====================
    \section{Mô hình hóa Logic và Các thuật toán Suy diễn}

    \subsection{Mô hình hóa bài toán bằng Logic Mệnh đề}

\subsubsection{Biến mệnh đề và Ngôn ngữ manh mối cốt lõi}
Trên bàn chơi kích thước $N \times N$, hệ thống biểu diễn nội bộ mỗi ô bởi cấu trúc \texttt{Cell(row, col)} với chỉ số bắt đầu từ $0$ ($0 \le \text{row}, \text{col} < N$) và ánh xạ hai chiều với tọa độ hiển thị dạng $\langle\text{Cột}\rangle\langle\text{Hàng}\rangle$ (cột đánh chữ cái từ \texttt{A}, hàng đánh số từ $1$, ví dụ \texttt{Cell(1, 2)} tương ứng với ô \texttt{C2} qua \texttt{parse\_cell} và \texttt{format\_cell} trong \texttt{griductive/regions.py}). Mỗi nhân vật $i$ trên lưới tương ứng với một \emph{biến mệnh đề chính} (primary variable) $C_i$:
\[
C_i =
\begin{cases}
\text{True},  & \text{nếu nhân vật } i \text{ là \emph{Criminal} (tội phạm)}, \\
\text{False}, & \text{nếu nhân vật } i \text{ là \emph{Innocent} (người vô tội)}.
\end{cases}
\]

Để bảo đảm tính tái lập giữa các lượt chạy và giữa tác tử suy diễn với bộ kiểm định, \texttt{BasicCNFEncoder} (\texttt{griductive/encoder.py}) khởi tạo ánh xạ biến \texttt{VarMap} một cách \emph{tất định}:
\begin{itemize}
    \item \textbf{Đánh chỉ số biến chính theo thứ tự \emph{row-major}:} Hàm \texttt{\_create\_var\_map} sắp xếp toàn bộ ô trong \texttt{board.characters} theo thứ tự từ điển $(\text{row}, \text{col})$ để tạo bộ \texttt{VarMap.primary} có độ dài $N^2$. Ô ở vị trí $idx \in \{0, \dots, N^2 - 1\}$ được gán định danh biến SAT cố định là:
    \[
    \text{id}(r, c) = idx + 1 = r \cdot N + c + 1 \in \{1, 2, \dots, N^2\}.
    \]
    Ví dụ, trên bàn $3 \times 3$, các ô \texttt{A1, B1, C1, A2, B2, C2, A3, B3, C3} lần lượt mang định danh $1, 2, \dots, 9$. Trong công thức CNF (\texttt{Literal = int} khác $0$), literal dương $+id$ biểu diễn $C_i$ (\texttt{CRIMINAL}) và literal âm $-id$ biểu diễn $\neg C_i$ (\texttt{INNOCENT}).

    \item \textbf{Tách biệt biến phụ:} Khi mã hóa các ràng buộc phức tạp (như manh mối đếm hoặc chẵn lẻ), bộ mã hóa cấp phát tuần tự các \emph{biến phụ} (auxiliary variables) thông qua hàm \texttt{allocate\_aux()} bắt đầu từ $N^2 + 1$ đến $N^2 + \text{auxiliary\_count}$, lưu tách biệt qua trường \texttt{VarMap.auxiliary\_count}. Mỗi manh mối được cấp phát tập biến phụ độc lập, không bao giờ dùng chung định danh giữa các manh mối khác nhau. Nhờ đó, tác tử suy diễn chỉ cần đặt giả định (\texttt{assumptions}) và phân loại trên tập biến chính $\{1, \dots, N^2\}$, đồng thời bộ kiểm tra tính duy nhất (\texttt{SATPuzzleValidator.check\_uniqueness}) chỉ dựng mệnh đề chặn nghiệm trên các biến chính mà không xem sai khác ở biến phụ là nghiệm mới.
\end{itemize}

Các manh mối có phạm vi vùng được chuẩn hóa qua \texttt{GridRegionResolver.resolve} (\texttt{griductive/regions.py}) để trả về một tuple các \texttt{Cell} phân biệt theo thứ tự \emph{row-major} (nếu tọa độ ngoài biên hoặc danh sách tường minh có phần tử trùng lặp, hàm ném ngoại lệ \texttt{ValueError}):
\begin{itemize}
    \item \texttt{RowRegion(row)} và \texttt{ColumnRegion(col)}: toàn bộ $N$ ô thuộc hàng \texttt{row} hoặc cột \texttt{col} (trong JSON đánh số hàng từ $1$ và cột theo chữ cái \texttt{A..Z}, mở rộng \texttt{AA..} theo hệ cơ số 26).
    \item \texttt{NeighborRegion(cell)}: các ô lân cận ngang, dọc, chéo quanh \texttt{cell} và loại trừ chính \texttt{cell}, không cuộn qua biên bàn (ô góc có $3$, ô biên có $5$, ô giữa có $8$ lân cận).
    \item \texttt{ExplicitRegion(cells)}: tập hợp tường minh các ô phân biệt (cho phép vùng rỗng).
\end{itemize}

Trong \texttt{griductive/models.py}, 6 mẫu manh mối cốt lõi được mô hình hóa bằng ba lớp bất biến \texttt{FactClue}, \texttt{BinaryClue} và \texttt{CountClue}. Lưu ý rằng mọi manh mối đã lật đều mang giá trị chân lý đúng bất kể ô sở hữu (\texttt{owner}) là \emph{Criminal} hay \emph{Innocent}. Quy ước $C_i \in \{0, 1\}$ khi lấy tổng (với $\text{True} \equiv 1$, $\text{False} \equiv 0$), ngữ nghĩa logic của 6 mẫu cốt lõi được tóm tắt trong Bảng~\ref{tab:core_clues}.

\begin{table}[H]
\centering
\footnotesize
\setlength{\tabcolsep}{4.5pt}
\renewcommand{\arraystretch}{1.25}
\begin{tabular}{@{}llllp{0.31\textwidth}@{}}
\toprule
\textbf{Mẫu} & \textbf{Lớp} & \textbf{Tham số} & \textbf{Ngữ nghĩa logic} & \textbf{Ý nghĩa và điều kiện} \\
\midrule
\texttt{FACT} & \texttt{FactClue} & \texttt{target}, \texttt{status} &
$C_t$ hoặc $\neg C_t$ &
Ô \texttt{target} ($t$) là \emph{Criminal} ($C_t$) hoặc \emph{Innocent} ($\neg C_t$). \\
\texttt{SAME} & \texttt{BinaryClue} & \texttt{left}, \texttt{right} &
$C_l \leftrightarrow C_r$ &
Hai nhân vật \texttt{left} ($l$) và \texttt{right} ($r$) có cùng trạng thái. \\
\texttt{DIFFERENT} & \texttt{BinaryClue} & \texttt{left}, \texttt{right} &
$C_l \oplus C_r$ &
Hai nhân vật \texttt{left} ($l$) và \texttt{right} ($r$) khác trạng thái nhau. \\
\texttt{EXACTLY} & \texttt{CountClue} & $k$, \texttt{region} ($R$) &
$\sum_{i \in R} C_i = k$ &
Có đúng $k$ tội phạm trong vùng $R$ ($k \in \mathbb{Z},\; 0 \le k \le |R|$); $k=0$ ứng với ``không có tội phạm''. \\
\texttt{AT\_LEAST} & \texttt{CountClue} & $k$, \texttt{region} ($R$) &
$\sum_{i \in R} C_i \ge k$ &
Có ít nhất $k$ tội phạm trong vùng $R$ ($k \in \mathbb{Z},\; 0 \le k \le |R|$). \\
\texttt{AT\_MOST} & \texttt{CountClue} & $k$, \texttt{region} ($R$) &
$\sum_{i \in R} C_i \le k$ &
Có tối đa $k$ tội phạm trong vùng $R$ ($k \in \mathbb{Z},\; 0 \le k \le |R|$). \\
\bottomrule
\end{tabular}
\caption{Ngữ nghĩa logic mệnh đề của 6 mẫu manh mối cốt lõi.}
\label{tab:core_clues}
\end{table}

\subsubsection{Các manh mối mở rộng (Clue Extensions)}
Bên cạnh 6 mẫu cốt lõi, nhóm đã thiết kế và hiện thực đầy đủ \textbf{2 mẫu manh mối mở rộng} khác biệt về cấu trúc lẫn phương pháp mã hóa, đồng bộ từ tầng mô hình (\texttt{griductive/models.py}), bộ đọc JSON (\texttt{griductive/puzzle\_io.py}), bộ đánh giá trực tiếp (\texttt{DirectClueEvaluator}), bộ mã hóa CNF (\texttt{BasicCNFEncoder}) đến hiển thị và tô sáng vùng trên giao diện (\texttt{griductive/gui/text.py}, \texttt{griductive/gui/highlight.py}), và đưa vào màn chơi thực nghiệm \texttt{4x4-04} ở Phần 5:

\begin{enumerate}
    \item \textbf{Manh mối chẵn lẻ --- \texttt{PARITY} (\texttt{ParityClue}):}
    \begin{itemize}
        \item \emph{Cú pháp:} Biểu diễn bởi \texttt{ParityClue(id, owner, kind, region)} với \texttt{kind} thuộc \texttt{ParityKind} (\texttt{EVEN} hoặc \texttt{ODD}). Trong định dạng JSON, manh mối khai báo \texttt{"type": "PARITY"}, trường \texttt{"parity"} nhận giá trị \texttt{"EVEN"} hoặc \texttt{"ODD"}, cùng một vùng \texttt{"region"} hợp lệ (ví dụ tại ô \texttt{B4} trong \texttt{4x4-04}: \texttt{PARITY(ODD, ROW 2)}).
        \item \emph{Ngữ nghĩa logic:} Với tập ô $R = \{i_1, i_2, \dots, i_m\}$ thu được từ vùng tham chiếu, manh mối ràng buộc tổng số \emph{Criminal} trong $R$ là số chẵn hoặc số lẻ:
        \[
        \text{PARITY}(\text{EVEN}, R) \equiv \left(\bigoplus_{t=1}^{m} C_{i_t} = \text{False}\right) \iff \left(\sum_{t=1}^{m} C_{i_t}\right) \bmod 2 = 0,
        \]
        \[
        \text{PARITY}(\text{ODD}, R) \equiv \left(\bigoplus_{t=1}^{m} C_{i_t} = \text{True}\right) \iff \left(\sum_{t=1}^{m} C_{i_t}\right) \bmod 2 = 1.
        \]
        \item \emph{Cách xử lý:} Trong \texttt{DirectClueEvaluator}, bộ đánh giá đếm số ô \texttt{CRIMINAL} trong $R$ và kiểm tra \texttt{count \% 2}. Trong \texttt{BasicCNFEncoder} (\texttt{\_parity\_encode}), thay vì liệt kê $2^{m-1}$ tổ hợp chẵn lẻ, hệ thống dùng \emph{chuỗi cổng XOR kiểu Tseitin} trên dãy literal $(x_1, \dots, x_m)$:
        \begin{itemize}
            \item Nếu $m = 0$ (vùng rỗng): \texttt{EVEN} luôn đúng nên trả về CNF rỗng \texttt{()}; \texttt{ODD} vô nghiệm nên trả về \texttt{((),)} (mệnh đề rỗng gây \texttt{UNSAT}).
            \item Nếu $m = 1$: trả về trực tiếp một mệnh đề đơn vị $(x_1,)$ cho \texttt{ODD} hoặc $(\neg x_1,)$ cho \texttt{EVEN}.
            \item Nếu $m \ge 2$: với $m > 2$, cấp phát $m - 2$ biến phụ $a_1, \dots, a_{m-2}$ để tính dần tổng XOR tiền tố $a_1 \leftrightarrow (x_1 \oplus x_2)$, $a_t \leftrightarrow (a_{t-1} \oplus x_{t+1})$. Mỗi cổng XOR $u \oplus v \leftrightarrow r$ được \texttt{\_xor\_encode} khai triển thành 4 mệnh đề 3-literal:
            \[
            (\neg u \lor \neg v \lor \neg r) \land (u \lor v \lor \neg r) \land (u \lor \neg v \lor r) \land (\neg u \lor v \lor r).
            \]
            Cuối cùng, đầu ra trung gian cuối $\mathit{cur}$ ($x_1$ nếu $m=2$, hoặc $a_{m-2}$ nếu $m \ge 3$) được ràng buộc với phần tử cuối $x_m$ bằng 2 mệnh đề 2-literal: $(\neg \mathit{cur} \lor x_m) \land (\mathit{cur} \lor \neg x_m)$ cho \texttt{EVEN} ($\mathit{cur} \leftrightarrow x_m$), hoặc $(\mathit{cur} \lor x_m) \land (\neg \mathit{cur} \lor \neg x_m)$ cho \texttt{ODD} ($\mathit{cur} \oplus x_m$). Tổng chi phí tăng tuyến tính theo $m$: đúng $m - 2$ biến phụ và $4m - 6$ mệnh đề với mọi $m \ge 2$.
        \end{itemize}
    \end{itemize}

    \item \textbf{Manh mối giao vùng --- \texttt{INTERSECTION} (\texttt{IntersectionClue}):}
    \begin{itemize}
        \item \emph{Cú pháp:} Biểu diễn bởi \texttt{IntersectionClue(id, owner, k, regions)} với $k \in \mathbb{Z}$ và \texttt{regions} là một tuple gồm các vùng. Trong tệp JSON, manh mối có \texttt{"type": "INTERSECTION"}, trường \texttt{"k"} và mảng \texttt{"regions"} chứa tối thiểu 2 vùng hợp lệ (ví dụ tại ô \texttt{C1} trong \texttt{4x4-04}: \texttt{INTERSECTION(k=0, [NEIGHBORS(B4), NEIGHBORS(B2)])}).
        \item \emph{Ngữ nghĩa logic:} Đặt $S = \bigcap_{j=1}^{p} R_j$ là giao của các tập ô đã giải từ \texttt{regions}. Manh mối khẳng định có đúng $k$ nhân vật \emph{Criminal} nằm đồng thời trong tất cả các vùng thành phần:
        \[
        \sum_{i \in S} C_i = k, \qquad 0 \le k \le |S|.
        \]
        Khi giao của các vùng là tập rỗng ($S = \emptyset$), manh mối vẫn hợp lệ nếu $k = 0$ (mã hóa thành CNF rỗng \texttt{()}) và ném \texttt{ValueError} nếu $k > |S|$.
        \item \emph{Cách xử lý:} Hàm \texttt{\_get\_intersection\_cells} (và \texttt{intersect\_regions}) giải từng vùng qua \texttt{GridRegionResolver}, lấy giao liên tiếp và sắp xếp kết quả theo thứ tự \emph{row-major} để cố định thứ tự biến. \texttt{DirectClueEvaluator} đếm số \texttt{CRIMINAL} trên tập $S$ và kiểm tra \texttt{count == k}. Hàm \texttt{\_intersection\_encode} của \texttt{BasicCNFEncoder} ánh xạ các ô trong $S$ thành literal dương rồi tái sử dụng hàm \texttt{\_encode\_count\_literals} với loại \texttt{CountKind.EXACTLY}, ngưỡng $k$ và bộ cấp phát \texttt{allocate\_aux}. Trên giao diện, \texttt{clue\_focus} tô sáng riêng các ô thuộc tập giao $S$ và tách các ô tâm của \texttt{NeighborRegion} vào nhóm neo (\texttt{anchors}), còn \texttt{describe\_clue} diễn đạt bằng cấu trúc \texttt{"both ... and ..."} để tránh nhầm lẫn giữa giao và hợp vùng.
    \end{itemize}
\end{enumerate}

    \subsection{Chuyển đổi tự động sang Dạng chuẩn Hội (CNF Encoding)}

\subsubsection{Cơ chế xây dựng cơ sở tri thức \texorpdfstring{$KB_t$}{KB\_t}}
Ở bước suy diễn $t$, phương thức \texttt{build\_kb(state)} của lớp \texttt{GriductiveLogicAgent} (trong \texttt{griductive/agent/inference.py}) chỉ tiếp nhận bản chụp công khai bất biến \texttt{PublicState(board, verdicts, clues)} từ engine, tuyệt đối không truy cập đáp án ẩn hay manh mối của các thẻ chưa lật. Tác tử chuyển toàn bộ thông tin công khai tại bước $t$ cho phương thức \texttt{encode} của \texttt{BasicCNFEncoder} để sinh tự động công thức CNF hoàn chỉnh:
\[
KB_t = \left(\bigwedge_{c \in \text{clues}_t} \text{CNF}(c)\right) \land \left(\bigwedge_{v \in \text{verdicts}_t} \text{Unit}(v)\right).
\]
Quy trình dựng $KB_t$ không sử dụng bất kỳ công thức CNF viết tay nào cho riêng từng màn chơi và diễn ra qua bốn bước:
\begin{enumerate}
    \item \textbf{Khởi tạo không gian biến:} Tạo \texttt{VarMap} gồm $N^2$ biến chính (định danh $1 \dots N^2$) theo thứ tự \emph{row-major} và khởi tạo bộ đếm \texttt{auxiliary\_count = 0}.
    \item \textbf{Chuyển đổi tự động từng manh mối công khai sang tập mệnh đề tuyển:} Duyệt tuần tự từng manh mối $c \in \text{clues}_t$ và gọi bộ mã hóa tương ứng:
    \begin{itemize}
        \item \texttt{FactClue} (\texttt{\_fact\_encode}): sinh $1$ mệnh đề đơn vị $((l,),)$.
        \item \texttt{BinaryClue} (\texttt{\_binary\_encode}): sinh trực tiếp $2$ mệnh đề $2$-literal biểu diễn tương đương hoặc XOR.
        \item \texttt{CountClue} (\texttt{\_count\_encode} $\to$ \texttt{\_encode\_count\_literals}): kiểm tra $k \in \mathbb{Z}$ và $0 \le k \le n$ với $n = |R|$, sau đó quy cả ba dạng đếm về \texttt{at\_most\_encode}:
        \begin{itemize}
            \item \texttt{AT\_MOST(k)}: mã hóa trực tiếp $\sum_{t=1}^{n} x_t \le k$ trên bộ literal dương $(x_1, \dots, x_n)$ theo \emph{bộ đếm tuần tự (Sequential Counter Encoding)} sử dụng phép định nghĩa tương đương hai chiều ($\leftrightarrow$) cho các biến trạng thái tổng tiền tố $s_{i,j}$ ($1 \le i \le n - 1,\; 1 \le j \le \min(i, k)$). Hai trường hợp biên được tối giản: nếu $k \ge n$ trả về \texttt{()} (luôn thỏa mãn); nếu $k = 0$ sinh $n$ mệnh đề đơn vị $(\neg x_1), \dots, (\neg x_n)$.
            \item \texttt{AT\_LEAST(k)}: áp dụng biến đổi đối ngẫu $\sum_{t=1}^{n} x_t \ge k \iff \sum_{t=1}^{n} (1 - x_t) \le n - k$, đảo dấu toàn bộ literal thành $(\neg x_1, \dots, \neg x_n)$ và gọi \texttt{at\_most\_encode} với ngưỡng $n - k$.
            \item \texttt{EXACTLY(k)}: kết hợp hội hai tập mệnh đề của \texttt{AT\_MOST(k)} trên $(x_1, \dots, x_n)$ và \texttt{AT\_MOST(n - k)} trên $(\neg x_1, \dots, \neg x_n)$, mỗi chiều sử dụng tập biến phụ riêng biệt.
        \end{itemize}
        \item \texttt{ParityClue} và \texttt{IntersectionClue}: lần lượt chuyển sang chuỗi XOR Tseitin hoặc đếm \texttt{EXACTLY(k)} trên tập giao $S$ như đã mô tả ở Mục 4.1.2.
    \end{itemize}
    \item \textbf{Biểu diễn phán quyết đã biết thành mệnh đề đơn vị:} Với mỗi phán quyết \texttt{Verdict(cell, status)} trong \texttt{state.verdicts} (bao gồm các thẻ mở ban đầu và những ô đã chứng minh ở các bước trước), hàm \texttt{\_get\_literal} ánh xạ sang literal $+id$ (nếu \texttt{CRIMINAL}) hoặc $-id$ (nếu \texttt{INNOCENT}) và nối thêm mệnh đề đơn vị $(l,)$ vào danh sách mệnh đề. Nếu dữ liệu đầu vào chứa mâu thuẫn logic, CNF vẫn được sinh hợp lệ để bộ giải DPLL phát hiện trạng thái \texttt{UNSAT}.
    \item \textbf{Đóng gói và thống kê kích thước $KB_t$:} Hàm trả về \texttt{EncodedKB(cnf, variables)} với \texttt{cnf} là \texttt{tuple[tuple[int, ...], ...]}. Từ đây hệ thống thống kê trực tiếp cho từng bước: số biến chính là $|\text{variables.primary}| = N^2$; số biến phụ là \texttt{variables.auxiliary\_count} (riêng mỗi lần gọi \texttt{at\_most\_encode(n, k)} với $1 \le k < n$ sinh đúng $\sum_{i=1}^{n-1}\min(i, k) = kn - \frac{k(k+1)}{2}$ biến phụ); và tổng số mệnh đề CNF là \texttt{len(kb.cnf)}. Khi chạy thực nghiệm toàn màn ở Phần 5, các chỉ số \texttt{peak\_auxiliary\_variables} và \texttt{peak\_clauses} lấy giá trị cực đại trên tất cả các $KB_t$ đã truy vấn trong vòng suy diễn.
\end{enumerate}

\subsubsection{Ví dụ dẫn xuất CNF cho manh mối nhị phân và manh mối đếm}
Để minh họa chi tiết quá trình chuyển đổi tự động trong \texttt{BasicCNFEncoder}, ta xét các manh mối thực tế trên bàn $3 \times 3$ của màn \texttt{3x3-01} (\texttt{puzzles/3x3-01.json}), với biến chính theo thứ tự \emph{row-major} là $C_1 \dots C_9$ ứng với \texttt{A1, B1, C1, A2, B2, C2, A3, B3, C3}.

\paragraph{1. Dẫn xuất cho manh mối nhị phân (\texttt{DIFFERENT} và \texttt{SAME}):}
Xét manh mối \texttt{clue-B1} có dạng \texttt{DIFFERENT(A1, B1)} và \texttt{clue-B3} có dạng \texttt{SAME(C2, C3)}:
\begin{itemize}
    \item Với \texttt{DIFFERENT(A1, B1)}: \texttt{A1} ứng với $C_1$ (ID $1$), \texttt{B1} ứng với $C_2$ (ID $2$). Biểu thức logic $C_1 \oplus C_2$ tương đương ``ít nhất một trong hai là \emph{Criminal} và không thể cả hai cùng là \emph{Criminal}'':
    \[
    C_1 \oplus C_2 \equiv (C_1 \lor C_2) \land \neg(C_1 \land C_2) \equiv (C_1 \lor C_2) \land (\neg C_1 \lor \neg C_2).
    \]
    Kết quả CNF thu được là \texttt{((1, 2), (-1, -2))} gồm $2$ mệnh đề và $0$ biến phụ.

    \item Với \texttt{SAME(C2, C3)}: \texttt{C2} ứng với $C_6$ (ID $6$), \texttt{C3} ứng với $C_9$ (ID $9$). Phép tương đương logic được phân rã thành hai phép kéo theo:
    \[
    C_6 \leftrightarrow C_9 \equiv (C_6 \to C_9) \land (C_9 \to C_6) \equiv (\neg C_6 \lor C_9) \land (C_6 \lor \neg C_9).
    \]
    Kết quả CNF là \texttt{((-6, 9), (6, -9))} gồm $2$ mệnh đề và $0$ biến phụ.
\end{itemize}

\paragraph{2. Dẫn xuất tổng quát cho bộ đếm tuần tự (\texttt{at\_most\_encode}):}
Với dãy $n$ literal $(x_1, \dots, x_n)$ và ngưỡng $1 \le k < n$, mỗi biến phụ $s_{i,j}$ ($1 \le i \le n - 1,\; 1 \le j \le \min(i, k)$) biểu diễn trạng thái ``trong $i$ phần tử đầu tiên $x_1, \dots, x_i$ có ít nhất $j$ literal nhận giá trị \text{True}''. Hàm \texttt{at\_most\_encode} chuyển các định nghĩa đệ quy sang CNF như sau:
\begin{itemize}
    \item \emph{Trường hợp cơ sở ($i = 1, j = 1$):} $s_{1,1} \leftrightarrow x_1 \equiv (\neg s_{1,1} \lor x_1) \land (s_{1,1} \lor \neg x_1)$.
    \item \emph{Biên dưới ($j = 1,\; i \ge 2$):} $s_{i,1} \leftrightarrow (s_{i-1,1} \lor x_i)$:
    \[
    (\neg s_{i,1} \lor s_{i-1,1} \lor x_i) \land (s_{i,1} \lor \neg s_{i-1,1}) \land (s_{i,1} \lor \neg x_i).
    \]
    \item \emph{Đường chéo ($2 \le j = i \le k$):} Do $s_{i-1,i}$ luôn sai nên $s_{i,i} \leftrightarrow (x_i \land s_{i-1,i-1})$:
    \[
    (\neg s_{i,i} \lor x_i) \land (\neg s_{i,i} \lor s_{i-1,i-1}) \land (s_{i,i} \lor \neg x_i \lor \neg s_{i-1,i-1}).
    \]
    \item \emph{Tổng quát ($2 \le j < i$):} Định nghĩa $s_{i,j} \leftrightarrow \big(s_{i-1,j} \lor (x_i \land s_{i-1,j-1})\big)$ được tách hai chiều và dùng luật phân phối của $\lor$ qua $\land$:
    \begin{align*}
    \text{Chiều thuận } (\to): &\quad \neg s_{i,j} \lor s_{i-1,j} \lor (x_i \land s_{i-1,j-1}) \\
    &\equiv (\neg s_{i,j} \lor s_{i-1,j} \lor x_i) \land (\neg s_{i,j} \lor s_{i-1,j} \lor s_{i-1,j-1}), \\
    \text{Chiều ngược } (\leftarrow): &\quad \neg\big(s_{i-1,j} \lor (x_i \land s_{i-1,j-1})\big) \lor s_{i,j} \\
    &\equiv (s_{i,j} \lor \neg s_{i-1,j}) \land (s_{i,j} \lor \neg x_i \lor \neg s_{i-1,j-1}).
    \end{align*}
    \item \emph{Ràng buộc cấm vượt quá $k$ ($k + 1 \le i \le n$):} Không được đồng thời $x_i$ đúng và đã có sẵn $k$ phần tử đúng trước đó:
    \[
    \neg(x_i \land s_{i-1,k}) \equiv (\neg x_i \lor \neg s_{i-1,k}).
    \]
\end{itemize}

\paragraph{3. Dẫn xuất cụ thể cho manh mối đếm \texttt{EXACTLY}:}
Xét manh mối mở đầu \texttt{clue-B2} của màn \texttt{3x3-01}: \texttt{EXACTLY(1, NEIGHBORS(A1))}. Vùng \texttt{NEIGHBORS(A1)} giải ra $3$ ô theo thứ tự \emph{row-major} là \texttt{B1} ($C_2$), \texttt{A2} ($C_4$), \texttt{B2} ($C_5$), tức $n = 3$, $k = 1$, bộ literal $(x_1, x_2, x_3) = (C_2, C_4, C_5) \equiv (2, 4, 5)$. Bộ mã hóa tách thành hội của hai ràng buộc:
\begin{enumerate}
    \item \textbf{Chiều \texttt{AT\_MOST(1)} trên $(2, 4, 5)$:} Cấp phát $1 \cdot 3 - \frac{1 \cdot 2}{2} = 2$ biến phụ là $s_{1,1}$ (ID $10$) và $s_{2,1}$ (ID $11$):
    \begin{itemize}
        \item $i=1, j=1$: $s_{1,1} \leftrightarrow C_2 \implies (\neg s_{1,1} \lor C_2) \land (s_{1,1} \lor \neg C_2)$, tức \texttt{(-10, 2), (10, -2)}.
        \item $i=2, j=1$: $s_{2,1} \leftrightarrow (s_{1,1} \lor C_4) \implies$ \texttt{(-11, 10, 4), (11, -10), (11, -4)}.
        \item Mệnh đề cấm cho $i=2, 3$: $\neg(C_4 \land s_{1,1})$ và $\neg(C_5 \land s_{2,1}) \implies$ \texttt{(-4, -10), (-5, -11)}.
    \end{itemize}
    Phần này sinh $2$ biến phụ (ID $10, 11$) và $7$ mệnh đề CNF.

    \item \textbf{Chiều \texttt{AT\_LEAST(1)} $\equiv$ \texttt{AT\_MOST(2)} trên $(-2, -4, -5)$:} Đặt $y_1 = \neg C_2$, $y_2 = \neg C_4$, $y_3 = \neg C_5$ với ngưỡng $k' = n - k = 2$. Cấp phát $2 \cdot 3 - \frac{2 \cdot 3}{2} = 3$ biến phụ mới là $s'_{1,1}$ (ID $12$), $s'_{2,1}$ (ID $13$) và $s'_{2,2}$ (ID $14$):
    \begin{itemize}
        \item $i=1, j=1$: $s'_{1,1} \leftrightarrow \neg C_2 \implies$ \texttt{(-12, -2), (12, 2)}.
        \item $i=2, j=1$: $s'_{2,1} \leftrightarrow (s'_{1,1} \lor \neg C_4) \implies$ \texttt{(-13, 12, -4), (13, -12), (13, 4)}.
        \item $i=2, j=2$ (nhánh đường chéo $j = i$): $s'_{2,2} \leftrightarrow (\neg C_4 \land s'_{1,1}) \implies$ \texttt{(-14, -4), (-14, 12), (14, 4, -12)}.
        \item Mệnh đề cấm cho $i=3$: $\neg(\neg C_5 \land s'_{2,2}) \equiv (C_5 \lor \neg s'_{2,2}) \implies$ \texttt{(5, -14)}.
    \end{itemize}
    Phần này sinh $3$ biến phụ (ID $12 \dots 14$) và $9$ mệnh đề CNF.
\end{enumerate}
Tổng hợp lại, riêng manh mối \texttt{clue-B2} tạo ra $5$ biến phụ và $16$ mệnh đề CNF. Khi ghép thêm phán quyết mở sẵn của thẻ \texttt{B2} là \texttt{CRIMINAL} thành mệnh đề đơn vị \texttt{(5,)}, $KB_0$ ban đầu của màn \texttt{3x3-01} có $9$ biến chính, $5$ biến phụ và $17$ mệnh đề; chỉ riêng lan truyền đơn vị từ \texttt{(5,)} qua \texttt{(-5, -11)}, \texttt{(11, -10)}, \texttt{(11, -4)} và \texttt{(10, -2)} đã lập tức suy ra $\neg C_4$ và $\neg C_2$, chứng minh được ngay \texttt{A2} và \texttt{B1} đều là \texttt{INNOCENT}.

    \subsection{Bộ giải DPLL SAT Solver}
\texttt{DPLLSolver} (\texttt{solver.py}) là bộ giải do nhóm tự cài đặt theo thủ tục DPLL~\cite{dpll1962}, không dùng thư viện SAT. Đầu vào gồm CNF (mỗi literal là số nguyên khác $0$, dấu dương nghĩa là \emph{Criminal}), bảng biến và một bộ giả thiết tạm thời. Đầu ra là \texttt{SAT} kèm giá trị cho mọi biến, kể cả biến phụ, hoặc \texttt{UNSAT}. Giả thiết chỉ có hiệu lực trong một lần gọi, nên tác tử kiểm tra được $KB_t \wedge C_i$ và $KB_t \wedge \neg C_i$ trên cùng một $KB_t$ mà không phải dựng lại.

\subsubsection{Thuật toán}
Mã giả dưới đây viết ở dạng đệ quy. Hình~\ref{fig:dpll_flow} là luồng của bản cài đặt: lan truyền đơn vị dùng kỹ thuật hai literal canh~\cite{chaff2001}, và lời gọi đệ quy được thay bằng một ngăn xếp nhánh. Mỗi phần tử ngăn xếp tương ứng một khung đệ quy, nên thứ tự duyệt và kết quả giống hệt mã giả.

\begin{lstlisting}[basicstyle=\ttfamily\footnotesize,frame=single,columns=fullflexible]
DPLL(F, A):                          # A: partial assignment
    A <- UNIT_PROPAGATE(F, A)
    if some clause is false under A:    return UNSAT
    x <- smallest-id unassigned variable
    if x is None:                       return SAT, A
    if DPLL(F, A + {x = False}) is SAT: return it
    return DPLL(F, A + {x = True})      # backtrack, then flip

SOLVE(F, assumptions):
    drop duplicate literals and tautologies; empty clause => UNSAT
    A <- assumptions + unit clauses of F   # conflict => UNSAT
    return DPLL(F, A)
\end{lstlisting}

\begin{figure}[H]
\centering
\begin{tikzpicture}
\node[box] (prep) at (0,0) {Chuẩn hóa CNF};
\node[decision] (empty) at (0,-2.2) {Mệnh đề\\rỗng?};
\node[terminal] (unsat) at (4.8,-2.2) {UNSAT};
\node[box] (init) at (0,-4.6) {Gán giả thiết và\\mệnh đề đơn vị};
\node[box] (prop) at (0,-6.8) {Lan truyền đơn vị\\(hai literal canh)};
\node[box] (flip) at (4.8,-6.8) {Quay lui,\\đảo sang True};
\node[decision] (conf) at (0,-9.0) {Mâu thuẫn?};
\node[decision] (more) at (4.8,-9.0) {Còn nhánh\\chưa đảo?};
\node[decision] (left) at (0,-11.4) {Còn biến\\chưa gán?};
\node[box] (pick) at (-4.6,-11.4) {Chọn biến id nhỏ nhất,\\gán False, lưu nhánh};
\node[terminal] (sat) at (0,-13.6) {SAT + mô hình đầy đủ};
\draw[arr] (prep) -- (empty);
\draw[arr] (empty) -- node[lbl, above] {có} (unsat);
\draw[arr] (empty) -- node[lbl, right] {không} (init);
\draw[arr] (init) -- (prop);
\draw[arr] (prop) -- (conf);
\draw[arr] (conf) -- node[lbl, above] {có} (more);
\draw[arr] (more) -- node[lbl, right] {có} (flip);
\draw[arr] (flip) -- (prop);
\draw[arr] (more.east) -- ++(1.1,0) |- node[lbl, right, pos=0.25] {không} (unsat.east);
\draw[arr] (conf) -- node[lbl, right] {không} (left);
\draw[arr] (left) -- node[lbl, above] {có} (pick);
\draw[arr] (pick) |- (prop);
\draw[arr] (left) -- node[lbl, right] {không} (sat);
\end{tikzpicture}
\caption{Luồng xử lý của \texttt{DPLLSolver}. Mâu thuẫn khi gán giả thiết và mệnh đề đơn vị cũng dẫn tới \texttt{UNSAT}.}
\label{fig:dpll_flow}
\end{figure}

\subsubsection{Thống kê và ví dụ}
\begin{itemize}
    \item \texttt{decisions}: số lần gán thử một giá trị cho biến phân nhánh, tính cả lần đảo giá trị.
    \item \texttt{propagations}: số biến được gán do lan truyền đơn vị, gồm cả mệnh đề đơn vị có sẵn; giả thiết không được tính.
    \item \texttt{backtracks}: số lần bỏ một nhánh thất bại và hoàn tác các phép gán của nó.
    \item \texttt{runtime\_s}: thời gian của lần gọi, đo bằng \texttt{time.perf\_counter}.
\end{itemize}

Với $F = (x_1 \vee x_2) \wedge (x_1 \vee \neg x_2)$, bộ giải chạy như sau:
\begin{enumerate}
    \item Chọn $x_1 = \text{False}$ (decision 1). Lan truyền ép $x_2 = \text{True}$ (propagation 1), khi đó mệnh đề $(x_1 \vee \neg x_2)$ bị sai: mâu thuẫn.
    \item Quay lui (backtrack 1) và đảo $x_1 = \text{True}$ (decision 2).
    \item $x_2$ còn tự do nên được chọn với giá trị $\text{False}$ (decision 3). Mọi mệnh đề đều thỏa.
\end{enumerate}
Kết quả là \texttt{SAT} với $x_1 = \text{True}$, $x_2 = \text{False}$ và thống kê (3, 1, 1); con số này khớp với kết quả chạy thật của bộ giải.

\texttt{tests/test\_solver.py} kiểm tra các trường hợp biên (CNF rỗng, mệnh đề rỗng, lan truyền dây chuyền, quay lui rồi đảo giá trị, giả thiết không rò sang lần gọi sau) và đối chiếu bộ giải với phép vét cạn trên 400 CNF sinh ngẫu nhiên.

    \subsection{Tác tử suy diễn (Deductive Logic Agent) và Kiểm tra tính duy nhất}

    \subsubsection{Phân loại hệ quả logic và Vòng lặp suy diễn}
    \textit{[Trình bày phương pháp chứng minh phản chứng để phân loại \texttt{CRIMINAL}, \texttt{INNOCENT}, \texttt{UNKNOWN}, \texttt{INCONSISTENT}; mô tả vòng lặp suy diễn cập nhật $KB_t$, xuất vết suy diễn (\emph{deduction trace}) và cơ chế hoạt động của nút \texttt{Hint}.]}

    \subsubsection{Kiểm tra tính duy nhất của lời giải (Uniqueness Check)}
\label{sec:uniqueness}
\texttt{SATPuzzleValidator} (\texttt{validation.py}) chạy khi engine nạp màn, trên \emph{toàn bộ} manh mối kể cả thẻ úp; kết quả không bao giờ được đưa cho tác tử. Gọi $F$ là CNF của toàn bộ manh mối. Đáp án lưu trong file không được thêm vào $F$, vì nếu thêm thì $F$ luôn chỉ còn một nghiệm. Sau khi giải $F$ được mô hình $M$, bộ kiểm định thêm \emph{mệnh đề chặn nghiệm} trên $N^2$ biến chính:
\[
    \beta_M = \bigvee_{i=1}^{N^2} \ell_i, \qquad
    \ell_i = \begin{cases} \neg C_i & \text{nếu } M(C_i) = \text{True},\\ C_i & \text{nếu } M(C_i) = \text{False}. \end{cases}
\]
$\beta_M$ chỉ sai tại đúng cách xếp của $M$, nên $F \wedge \beta_M$ còn thỏa được khi và chỉ khi tồn tại cách xếp \emph{Criminal}/\emph{Innocent} thứ hai. Mệnh đề chặn không chứa biến phụ, vì hai mô hình chỉ khác nhau ở biến phụ vẫn là cùng một nghiệm. Do đó phép kiểm tra chỉ tốn đúng hai lượt gọi SAT.

\begin{figure}[H]
\centering
\begin{tikzpicture}
\node[terminal] (err) at (1.85,2.1) {Báo lỗi\\(không nạp)};
\node[box] (struct) at (0,0) {Cấu trúc\\hợp lệ?};
\node[box] (eval) at (3.7,0) {Đáp án thỏa\\mọi manh mối?};
\node[box] (solve1) at (7.3,0) {Giải $F$};
\node[decision] (d1) at (10.3,0) {SAT?};
\node[terminal] (unsat) at (13.4,0) {UNSAT};
\node[box] (solve2) at (10.3,-2.9) {Giải $F \wedge \beta_M$};
\node[decision] (d2) at (6.5,-2.9) {SAT?};
\node[terminal] (unique) at (2.6,-2.9) {UNIQUE\\(nạp màn)};
\node[terminal] (multi) at (6.5,-5.1) {MULTIPLE};
\draw[arr] (struct) -- node[lbl, above] {có} (eval);
\draw[arr] (eval) -- node[lbl, above] {có} (solve1);
\draw[arr] (solve1) -- (d1);
\draw[arr] (d1) -- node[lbl, above] {không} (unsat);
\draw[arr] (d1) -- node[lbl, right] {có} (solve2);
\draw[arr] (solve2) -- (d2);
\draw[arr] (d2) -- node[lbl, above] {không} (unique);
\draw[arr] (d2) -- node[lbl, right] {có} (multi);
\draw[arr] (struct.north) -- node[lbl, left] {không} (err);
\draw[arr] (eval.north) -- node[lbl, right] {không} (err);
\end{tikzpicture}
\caption{Luồng kiểm định khi nạp màn. Bước kiểm tra đáp án dùng bộ đánh giá trực tiếp, độc lập với CNF. Engine chỉ nạp màn có kết quả \texttt{UNIQUE}; mọi trường hợp khác giữ nguyên ván đang chơi.}
\label{fig:validation_flow}
\end{figure}

\texttt{tests/test\_validation.py} kiểm thử các trường hợp một nghiệm, hai nghiệm, manh mối mâu thuẫn, mô hình chỉ khác ở biến phụ và đáp án vi phạm manh mối, đồng thời chạy validator trên mọi file trong \texttt{puzzles/}.

\paragraph{Nghiệm duy nhất chưa đủ để giải mà không phải đoán:} Tính duy nhất xét toàn bộ manh mối cùng lúc, còn tác tử chỉ thấy các manh mối được lật dần. Chín màn của phiên bản dữ liệu đầu tiên đều \texttt{UNIQUE}, nhưng bốn màn khiến tác tử dừng ở \texttt{STUCK} (Bảng~\ref{tab:legacy_puzzles}). Vì vậy bộ màn trong \texttt{puzzles/} được thiết kế lại và phải thỏa thêm điều kiện ``tác tử giải trọn từ các thẻ mở ban đầu''.

    \begin{table}[H]
    \centering\small
    \begin{tabular}{@{}lccc@{}}
    \toprule
    \textbf{Màn (bản đầu)} & \textbf{Tính duy nhất} & \textbf{Kết quả tác tử} & \textbf{Số thẻ đã mở} \\
    \midrule
    \texttt{3x3-01} & \texttt{UNIQUE} & \texttt{SOLVED} & 9/9 \\
    \texttt{3x3-02} & \texttt{UNIQUE} & \texttt{STUCK} & 1/9 \\
    \texttt{3x3-03} & \texttt{UNIQUE} & \texttt{STUCK} & 1/9 \\
    \texttt{4x4-01}, \texttt{4x4-02}, \texttt{4x4-03} & \texttt{UNIQUE} & \texttt{SOLVED} & 16/16 \\
    \texttt{5x5-01} & \texttt{UNIQUE} & \texttt{SOLVED} & 25/25 \\
    \texttt{5x5-02} & \texttt{UNIQUE} & \texttt{STUCK} & 4/25 \\
    \texttt{5x5-03} & \texttt{UNIQUE} & \texttt{STUCK} & 2/25 \\
    \bottomrule
    \end{tabular}
    \caption{Chín màn của phiên bản dữ liệu đầu tiên: nghiệm duy nhất nhưng không phải màn nào cũng giải được mà không đoán.}
    \label{tab:legacy_puzzles}
    \end{table}


    % ==================== PHẦN 5 ====================
    \clearpage
    \section{Thực nghiệm và Đánh giá}
    \subsection{Thiết lập phép đo}
    Bộ thử gồm 10 màn chính thức trong \texttt{puzzles/} tại commit \texttt{0ee67a6}, cũng là các màn trong danh sách Load của game. Có ba màn $3\times3$, bốn màn $4\times4$ và ba màn $5\times5$. Mỗi màn chạy một lần trong tiến trình riêng trên Linux x86-64, CPU AMD Ryzen 7 6800H, Python 3.11.15, dùng cùng bộ mã hóa CNF và bộ giải DPLL. Lệnh tái lập là \verb|python main.py --benchmark ket-qua.csv --timeout 30|. File \texttt{report/benchmark.csv} lưu số liệu gốc.

    Đồng hồ bắt đầu sau khi nạp và kiểm định puzzle, dừng ngay khi vòng suy diễn trả về. Mỗi bước tính cả truy vấn SAT để Agent chọn nước đi và truy vấn SAT để Engine chấm nước đó. Một bước được tính khi Engine chấp nhận mở thẻ. Số biến phụ và số mệnh đề là giá trị lớn nhất trong các KB đã truy vấn. Giới hạn suy diễn là 30 giây cho mỗi màn. Cả 10 màn đều đạt \texttt{SOLVED} trong thời hạn.

    \subsection{Kết quả theo kích thước}
    \begin{table}[H]
    \centering\small
    \begin{tabular}{@{}lrrrrrr@{}}
    \toprule
    Bàn & Giải hết & Biến phụ & Mệnh đề & Lượt SAT & Bước & Thời gian (ms) \\
    \midrule
    $3 \times 3$ & 3/3 & 24.0 & 94.0 & 127.3 & 7.0 & 5.89 \\
    $4 \times 4$ & 4/4 & 95.2 & 362.5 & 448.0 & 14.0 & 49.53 \\
    $5 \times 5$ & 3/3 & 261.7 & 977.7 & 1028.7 & 21.7 & 329.17 \\
    \bottomrule
    \end{tabular}
    \caption{Số màn giải hết và giá trị trung bình mỗi màn theo kích thước bàn.}
    \label{tab:benchmark_summary}
    \end{table}
    \subsection{Kết quả từng màn}
    \begin{table}[H]
    \centering\small
    \setlength{\tabcolsep}{4pt}
    \begin{tabular}{@{}lrrrrrrrr@{}}
    \toprule
    Màn & Biến C/P & Mệnh đề & Lượt SAT & Chọn & Lan truyền & Quay lui & Bước & ms \\
    \midrule
    \texttt{3x3-01} & 9/19 & 78 & 160 & 0 & 2274 & 0 & 8 & 6.89 \\
    \texttt{3x3-02} & 9/28 & 111 & 126 & 0 & 1546 & 0 & 7 & 5.28 \\
    \texttt{3x3-03} & 9/25 & 93 & 96 & 0 & 1770 & 0 & 6 & 5.49 \\
    \texttt{4x4-01} & 16/105 & 385 & 448 & 242 & 18228 & 24 & 14 & 45.16 \\
    \texttt{4x4-02} & 16/125 & 471 & 448 & 768 & 26256 & 0 & 14 & 73.06 \\
    \texttt{4x4-03} & 16/102 & 386 & 448 & 148 & 17458 & 0 & 14 & 47.29 \\
    \texttt{4x4-04} & 16/49 & 208 & 448 & 210 & 11194 & 6 & 14 & 32.62 \\
    \texttt{5x5-01} & 25/179 & 666 & 1150 & 916 & 86626 & 0 & 23 & 233.54 \\
    \texttt{5x5-02} & 25/332 & 1246 & 880 & 1496 & 145006 & 84 & 20 & 394.12 \\
    \texttt{5x5-03} & 25/274 & 1021 & 1056 & 4160 & 137682 & 580 & 22 & 359.86 \\
    \bottomrule
    \end{tabular}
    \caption{C/P là số biến chính/phụ.}
    \label{tab:benchmark_detailed}
    \end{table}

    \subsection{Thời gian suy diễn}
    \begin{figure}[H]
    \centering
    \begin{tikzpicture}
    \begin{axis}[
        width=0.97\linewidth, height=5.3cm,
        ybar, bar width=12pt, ymin=0, ymax=450, ytick distance=100,
        ylabel={Thời gian (ms)}, ymajorgrids, axis lines=left,
        xtick=data, xticklabels={3x3-01,3x3-02,3x3-03,4x4-01,4x4-02,4x4-03,4x4-04,5x5-01,5x5-02,5x5-03},
        xticklabel style={font=\scriptsize}, enlarge x limits=0.08,
    ]
    \addplot[draw=blue!65!black,fill=blue!45] table[
        col sep=comma, x expr=\coordindex, y expr=1000*\thisrow{runtime_s}
    ] {benchmark.csv};
    \end{axis}
    \end{tikzpicture}
    \caption{Thời gian suy diễn của từng file trong \texttt{puzzles/}.}
    \end{figure}

    \subsection{Nhận xét}
    Trên cả ba màn $3\times3$, lan truyền đơn vị và kiểm tra mâu thuẫn đủ để trả lời mọi truy vấn SAT. Cả bốn màn $4\times4$ đều cần 14 bước và 448 lượt gọi SAT. Màn mới \texttt{4x4-04} mở sẵn hai ô C1 và C2, có cả manh mối \texttt{INTERSECTION} và \texttt{PARITY}. Bộ giải mở 14 ô còn lại trong 32.62 ms, nhanh nhất nhóm. Riêng màn \texttt{4x4-02} có KB đạt 471 mệnh đề và ghi nhận 26\,256 lần lan truyền, cao nhất nhóm. Thời gian đo được cũng dài nhất nhóm.

    Trên bàn $5\times5$, màn \texttt{5x5-02} chỉ cần 20 bước nhưng mất nhiều thời gian nhất trong bộ thử. Màn này có 1\,246 mệnh đề ở bước lớn nhất và 145\,006 lần lan truyền. Màn \texttt{5x5-03} có nhiều lần quay lui nhất (580). Chỉ đếm số thẻ mở không phản ánh được chi phí của từng truy vấn SAT. Thời gian trong bảng là một lần đo trên một máy; mức dao động giữa các lần chạy chưa được khảo sát.

    \subsection{Trường hợp thất bại và hạn chế}
    \paragraph{Trường hợp thất bại:} Trên bộ 10 màn chính thức, không có màn nào bị \texttt{TIMEOUT}, \texttt{ERROR}, \texttt{STUCK} hay \texttt{INCONSISTENT}. Tuy vậy, bộ màn này đã được thiết kế để giải được. Trên phiên bản dữ liệu đầu tiên (Bảng~\ref{tab:legacy_puzzles}), tác tử dừng ở \texttt{STUCK} tại bốn trên chín màn, dù cả chín màn đều có nghiệm duy nhất. Ở các màn này, tác tử không đoán mà trả về \texttt{STUCK} và giữ nguyên các thẻ còn úp, đúng như yêu cầu của đề.

    \paragraph{Hạn chế:}
    \begin{itemize}
        \item Bộ giải DPLL dùng quay lui theo thứ tự thời gian, chưa có học mệnh đề từ xung đột (CDCL) hay heuristic chọn biến động. Số lần chọn nhánh tăng từ $0$ ở bàn $3\times3$ lên khoảng 900 đến hơn 4\,000 ở bàn $5\times5$, kéo theo thời gian tăng nhanh (Bảng~\ref{tab:benchmark_detailed}).
        \item Mỗi bước, tác tử gọi bộ giải $1 + 2m$ lần với $m$ ô chưa biết, và engine lặp lại phép phân loại này khi chấm phán quyết. Hai phía đều đúng thiết kế tách vai trò nhưng làm số lượt gọi SAT tăng gấp đôi.
        \item Nút \texttt{Hint} chỉ ra manh mối đã lật đầu tiên có nhắc tới ô bị ép buộc. Manh mối này liên quan đến ô đó nhưng chưa chắc là manh mối thực sự quyết định phép chứng minh.
    \end{itemize}

    % ==================== PHẦN 6 ====================
    \section{Video Demo}
    \textit{[Đính kèm liên kết Google Drive công khai của video demo (dài 5--10 phút, có thuyết minh hoặc phụ đề, đã kiểm tra quyền xem bằng trình duyệt ẩn danh) và liệt kê ngắn gọn các nội dung trình diễn giao diện, chơi thủ công và chạy bộ giải tự động.]}

    % ==================== TÀI LIỆU THAM KHẢO ====================
    \begin{thebibliography}{9}
    \bibitem{griductive_game}
    Knotwise Games, \emph{Griductive}. Trang chơi: \url{https://knotwise.games/games/griductive}; hướng dẫn chơi: \url{https://knotwise.games/games/griductive/how-to-play}.

    \bibitem{russell2020aima}
    Stuart Russell and Peter Norvig, \emph{Artificial Intelligence: A Modern Approach}, 4th ed. Pearson, 2020.

    \bibitem{dpll1962}
    Martin Davis, George Logemann, and Donald Loveland, ``A Machine Program for Theorem-Proving,'' \emph{Communications of the ACM}, vol. 5, no. 7, pp. 394--397, 1962.

    \bibitem{chaff2001}
    Matthew W. Moskewicz, Conor F. Madigan, Ying Zhao, Lintao Zhang, and Sharad Malik, ``Chaff: Engineering an Efficient SAT Solver,'' in \emph{Proceedings of the 38th Design Automation Conference (DAC)}, pp. 530--535, 2001.

    \bibitem{pygame}
    Pygame Community, \emph{Pygame 2.6 Documentation}. \url{https://www.pygame.org/docs/}

    \bibitem{dragonball_api}
    Dragon Ball API, \emph{Characters} (ảnh nhân vật). \url{https://dragonball-api.com/}. Nhân vật \emph{Dragon Ball} \copyright{} Akira Toriyama / Bird Studio, Shueisha, Toei Animation.
    \end{thebibliography}

    \appendix
    \section{Phụ lục: Sử dụng công cụ Trí tuệ Nhân tạo Tạo sinh (Generative AI)}
% TODO: mỗi thành viên bổ sung công cụ AI đã dùng, mục đích và prompt tiêu biểu (ghi đúng thực tế).
Nhóm sử dụng các mô hình ngôn ngữ lớn (LLM) như công cụ hỗ trợ. Bộ mã hóa CNF, bộ giải DPLL và tác tử suy diễn đều do nhóm tự xây dựng, không dùng thư viện SAT có sẵn. Mọi nội dung do AI gợi ý đều được nhóm kiểm tra lại bằng mã nguồn, kiểm thử tự động và kết quả chạy thực tế. Bảng~\ref{tab:genai} tóm tắt cách sử dụng.

\begin{table}[H]
\centering
\small
\renewcommand{\arraystretch}{1.3}
\begin{tabular}{@{}p{0.14\textwidth} p{0.38\textwidth} p{0.42\textwidth}@{}}
\toprule
\textbf{Công cụ} & \textbf{Mục đích sử dụng} & \textbf{Prompt tiêu biểu} \\
\midrule
Claude Code & Đề xuất ý tưởng và chủ đề cho giao diện, đánh giá nguồn hình ảnh có thể dùng. Nhóm chốt chủ đề Dragon Ball. & ``Phần chốt GUI \ldots\ ý tưởng lên cho giao diện này sẽ như thế nào, vì sao lại lên như vậy, nếu lên như vậy thì liệu có nguồn hay không, và nếu không thì tạo bằng cách gì? Cho tôi một số chủ đề \ldots\ và nên có chủ đề Dragon Ball nha.'' \\
Claude Code & Thiết kế lại giao diện theo phong cách tối giản, tông tối, tham khảo giao diện gốc của Knotwise Games; thay ảnh nhân vật tự vẽ bằng ảnh nhân vật Dragon Ball. & ``Tôi đang muốn giao diện minimalism với chủ đề Dragon Ball, đồng thời tôi muốn ảnh thẻ nhân vật là các nhân vật trong Dragon Ball \ldots\ giao diện tham khảo ở Knotwise Games.'' và ``Tôi đang muốn làm theo theme tối màu, phải minimalism nữa.'' \\
Claude Code & Rà soát mức độ đáp ứng yêu cầu của đề, chụp ảnh giao diện và soạn nháp các mục báo cáo. & ``Giả sử bạn là chuyên gia trong lĩnh vực kiểm thử và đánh giá game, đọc nội dung đồ án trong thư mục data và kiểm lại xem đồ án của tôi đã đáp ứng đủ yêu cầu chưa.'' \\
% TODO: thêm dòng cho các thành viên khác, ví dụ:
% ChatGPT & Tìm hiểu ... & ``...'' \\
\bottomrule
\end{tabular}
\caption{Các công cụ AI tạo sinh đã sử dụng.}
\label{tab:genai}
\end{table}

\end{document}
